\documentclass[conference]{IEEEtran}
\IEEEoverridecommandlockouts
\usepackage{cite}
\usepackage{amsmath,amssymb,amsfonts}
\usepackage{graphicx}
\usepackage{textcomp}
\usepackage{xcolor}
\usepackage{booktabs}
\usepackage[hidelinks]{hyperref}
\def\BibTeX{{\rm B\kern-.05em{\sc i\kern-.025em b}\kern-.08em T\kern-.1667em\lower.7ex\hbox{E}\kern-.125emX}}

\begin{document}

\title{Knowing When You Do Not Know: Sequential Map-Anchored\\
Visual Localization for GNSS-Denied UAV Flight}

\author{\IEEEauthorblockN{Yusuf G\"unel}
\IEEEauthorblockA{\textit{Independent Researcher} \\
Ankara, T\"urkiye \\
yusufgunel71@hotmail.com}}

\maketitle

\begin{abstract}
Matching a downward-looking camera against a pre-loaded satellite map is the
standard route to absolute position when GNSS is denied. The recent literature
optimises this as a single-image retrieval problem, and the common benchmark
is close to saturated. A flying aircraft, however, produces a \emph{sequence},
and in that setting the binding constraint is not retrieval accuracy but the
fact that a large fraction of frames cannot be matched at all. We fuse
cross-view matching with visual odometry in a particle filter and evaluate on
ten real UAV survey flights spanning 406--2572\,m altitude, 9--103\,km per
flight and acquisition dates from 2016 to 2023. On the seven flights whose
imagery is usable against the basemap the system holds a median error of
8.4--24.8\,m with $\geq$99.7\% of frames localised, while visual odometry
alone drifts 2803\,m over 74\,km. Two by-products follow from the same
machinery: the residual rotation of each map match exposes the platform's
heading error, letting the system calibrate its own magnetometer deviation
without GNSS and cutting odometry drift from 3.795\% to 0.865\% of distance
travelled; and the \emph{match rate}, measurable on a planned route before
flight, predicts whether the system will work there at all
($\rho=-0.764$ against log error, with a threshold near 50\%). Finally we
report a negative result that we believe generalises: replacing LoFTR with
RoMa raises match rate on the failing flights from 12--42\% to 96--100\% and
makes end-to-end accuracy dramatically \emph{worse}, because RoMa returns
confident correspondences for unrelated imagery. For localization, a matcher
is useful in proportion to its ability to stay silent, not to match.
\end{abstract}

\begin{IEEEkeywords}
GNSS-denied navigation, cross-view geo-localization, particle filter, visual
odometry, unmanned aerial vehicles, image matching
\end{IEEEkeywords}

\section{Introduction}

An unmanned aerial vehicle whose GNSS receiver is jammed or spoofed loses
absolute position. Inertial and visual odometry continue to supply relative
motion, but their error accumulates without bound: on the 74\,km flight
studied here, frame-to-frame visual odometry ends 2803\,m from truth. The
established remedy is to carry a georeferenced satellite basemap and recover
absolute position by matching what the camera sees against it.

Recent work treats this as cross-view image retrieval and has driven the
common benchmark close to saturation; University-1652 now sits at 97.5\%
Recall@1~\cite{mgs2net}. That framing quietly assumes the hard part is
choosing correctly among candidate locations. In flight over real terrain we
find a different constraint dominates: on a substantial fraction of frames
\emph{no} correct match exists to be chosen. Over water, uniform farmland, or
repeated building patterns, and wherever the basemap and the camera disagree
in season or in years, matching returns nothing usable. In our single-frame
baseline this is 23\% of frames on the easiest flight in our set and 96\% on
the hardest.

A flying aircraft is not a collection of independent photographs. It follows a
trajectory, and frames that cannot be matched are bracketed by frames that
can. We therefore treat the problem as sequential state estimation: visual
odometry supplies a motion model, map matching supplies absolute fixes and is
permitted to fail, and a particle filter carries the two together. This is
standard practice in ground robotics but is almost absent from the recent UAV
cross-view literature, which we survey in Section~\ref{sec:related}.

Our contributions are:

\begin{enumerate}
\item A sequential map-anchored localization system evaluated on \emph{ten}
real UAV flights rather than one, spanning a 6.3$\times$ altitude range and
seven years of acquisition dates, fully automatic with per-flight calibration
derived from each flight's own first 20\% and evaluated on the remainder
(Section~\ref{sec:results}).
\item Online self-calibration of platform heading and image scale from the
geometry of the map matches themselves, requiring no GNSS. The measured
magnetometer deviation, $-1.93^\circ$ and $-7.08^\circ$ on the two legs of a
survey pattern, agrees with the GNSS-derived value in magnitude and sign, and
correcting for it reduces odometry drift by a factor of 4.4
(Section~\ref{sec:selfcal}).
\item A pre-flight predictor of deployability. The fraction of frames that
match the basemap at a known position separates working from failing flights
with $\rho = -0.764$ against log error and a threshold near 50\%. Because it
requires no flight, it turns "will this work here?" into a measurement
(Section~\ref{sec:law}).
\item A negative result on matcher choice, with the mechanism identified: a
stronger dense matcher degrades localization because it lacks the ability to
refuse (Section~\ref{sec:matcher}).
\end{enumerate}

\section{Related Work}
\label{sec:related}

\textbf{Cross-view geo-localization.} The dominant formulation matches a
single UAV image against a database of satellite tiles. Recent entries include
transformer and state-space architectures~\cite{dinogfsa,mgs2net}, robustness
to off-nadir viewing~\cite{offnadir}, degraded imagery~\cite{relate}, and
metric scale recovery via semantic anchors~\cite{scaleaware}. Surveying
twenty such papers from 2024--2026 we found none that fuse the per-frame
result with a motion model over time.

\textbf{Sequential localization against overhead imagery.} Particle filtering
against aerial maps is established for \emph{ground} platforms: BEV-Patch-PF
combines learned bird's-eye matching with a particle filter for off-road
vehicles~\cite{bevpatchpf}, and earlier work localises street-level imagery
city-wide to roughly 20\,m~\cite{downes}. For aircraft the closest work
predates modern learned matching: Shan et al.~\cite{shan} filter optical flow
against map correlation. Our contribution sits at the intersection: a modern
detector-free matcher inside a sequential filter, on an aircraft, evaluated
with metric error over multi-kilometre trajectories.

\textbf{Image matching.} LoFTR~\cite{loftr} performs detector-free coarse-to-fine
matching with transformers. RoMa~\cite{roma} predicts a dense warp with a
foundation-model encoder and leads current matching benchmarks. We use both
and report in Section~\ref{sec:matcher} why benchmark ranking does not transfer
to this task.

\section{Method}

\subsection{Frame preparation}
Each UAV frame is rotated to north-up using the platform heading and resampled
to the basemap's ground scale, giving a metric, north-aligned crop covering
$300 \times 300$\,m. Heading comes from the inertial unit and altitude from the
barometric altimeter; neither depends on GNSS, matching the pose-prior protocol
used by contemporary benchmarks. The basemap is stored in EPSG:4326, whose
pixels are square in degrees but not in metres --- at our latitude
0.2526\,m/px east-west against 0.2974\,m/px north-south, an 17.8\% anisotropy
--- so satellite crops are resampled to be metrically square before matching.

\subsection{Absolute fixes}
A grid of overlapping 300\,m basemap tiles is indexed with frozen DINOv2
descriptors~\cite{dinov2}. For a query frame, the nearest tiles by cosine
similarity give candidate regions; each candidate is verified by LoFTR matching
against a 400\,m crop, followed by MAGSAC homography estimation. The image
centre transported through the homography gives a ground point.

\subsection{From image centre to aircraft position}
\label{sec:boresight}
The matched ground point is where the optical axis meets the ground, not where
the aircraft is. With altitude $h$ and attitude angle $\theta$ the two differ
by $h\tan\theta$; at 466\,m a $2^\circ$ tilt displaces the estimate by 16\,m.
Regressing localization error in the body frame against attitude recovers this
relation with coefficients $+0.981$ (along-track) and $-0.972$ (cross-track),
confirming the model. Calibrating the constant mounting offset on a flight's
first 20\% and evaluating on the remaining 80\% reduces median error from
16.82\,m to 6.12\,m.

Two safeguards proved necessary across flights. First, the recorded height is
not always true height above ground; on one flight (recorded 2313\,m) applying
the correction \emph{increased} error from 27.8\,m to 131.2\,m, so an effective
altitude is estimated from the data. Second, calibration frames are split in
half, the correction is fitted on one half and verified on the other, and a
correction that does not reduce error is not applied. This declined to fire on
two of nine flights.

\subsection{Sequential fusion}
State is planar position. The motion model applies the odometry increment with
process noise matched to its measured step accuracy. The measurement model
treats candidate positions as a multi-modal likelihood with a flat outlier
floor. A Kalman filter is inappropriate here: map matching is usually accurate
to metres but occasionally confident and completely wrong, giving a
heavy-tailed, multi-modal error distribution that a single Gaussian mode cannot
represent.

Two mechanisms proved essential in practice. \emph{Measurement-driven
injection}: once the particle cloud contracts to a few metres, no particle lies
near a fix 100\,m away, the likelihood is uniformly zero and the filter cannot
move even when the measurement is correct --- we observed a filter locked
105\,m from truth for twenty frames while receiving 4--10\,m measurements.
Replacing a fraction of particles with draws from the measurement distribution
each step resolves this. \emph{Dead reckoning on the heading}: when odometry
fails, propagating the last motion \emph{vector} fails at turns; propagating
speed along the inertially-known heading does not.

Because the predicted position is known, verification normally requires a
single match attempt rather than a search over the map: 1.52 calls per frame
against 3.80 for the single-frame baseline.

\subsection{Self-calibration from the map}
\label{sec:selfcal}
The homography aligning a north-up frame to a north-up tile carries a residual
rotation. If heading were exact this residual would be zero; it is therefore a
direct measurement of heading error. Measured across a survey flight it was
$-1.93^\circ$ ($\sigma=0.86$) on the northwest legs and $-7.08^\circ$
($\sigma=1.73$) on the southeast legs. Heading error varying with heading is
the signature of magnetometer hard-iron distortion. The GNSS-derived odometry
direction bias on the same legs was $+1.65^\circ$ and $+7.22^\circ$: same
magnitudes, opposite sign, as the geometry predicts. Feeding a running estimate
of this residual back into the odometry reduced drift from 3.795\% to 0.865\%
of distance travelled, that is 2803\,m to 639\,m over 74\,km, using no
positioning information of any kind. The scale component of the same homography
similarly corrects the assumed ground sampling distance as terrain elevation
changes.

\section{Experimental Setup}

We use UAV-VisLoc~\cite{uavvisloc}, which pairs real survey flights with
georeferenced basemaps and post-processed GNSS. We verified the reference
quality independently: on straight survey legs the cross-track scatter is
1.50\,m and the step-length scatter 0.60\,m, so the reference is clean relative
to the errors we report.

Of eleven flights we evaluate ten. The excluded one has no attitude or heading
columns in its metadata, which the method requires; our loader rejects it
explicitly rather than producing unsupported numbers. One flight ships its
basemap as a $2\times2$ grid of GeoTIFFs totalling 4.5\,GB rather than a single
file; we present these to the pipeline through a GDAL virtual raster so that
no copy is made, which is also how tiled basemaps would arrive in deployment. Per-flight ground sampling distance and
camera boresight are estimated from each flight's own first 20\% and all
reported figures come from the remaining 80\%.

We report coverage (fraction of frames for which a position is produced),
median and 90th-percentile error, and matching calls per frame. All experiments
run on a laptop NVIDIA RTX 3050 Ti with 4\,GB of VRAM; peak usage is 1.4\,GB
with LoFTR.

\section{Results}
\label{sec:results}

\begin{table*}[t]
\caption{Per-flight results, fully automatic. The seven flights above the
rule have a match rate of at least 50\%; the three below do not.}
\label{tab:flights}
\centering
\begin{tabular}{lrrrrrrr}
\toprule
Flight & Frames & Altitude (m) & Distance (km) & Match rate & Coverage &
Median error (m) & 90th pct.\ (m) \\
\midrule
03 & 768  & 466  & 74  & 93\% & 100.0\% & \textbf{8.35}  & 20.03 \\
09 & 766  & 546  & 77  & 70\% & 100.0\% & \textbf{14.94} & 59.33 \\
06 & 344  & 834  & 24  & 76\% & 99.7\%  & \textbf{15.06} & 365.32 \\
04 & 738  & 544  & 83  & 90\% & 100.0\% & \textbf{15.44} & 54.64 \\
05 & 473  & 2313 & 30  & 50\% & 99.8\%  & \textbf{16.69} & 182.51 \\
01 & 817  & 406  & 66  & 77\% & 100.0\% & \textbf{22.51} & 114.60 \\
11 & 590  & 2572 & 84  & 90\% & 99.8\%  & \textbf{24.79} & 424.41 \\
\midrule
02 & 1071 & 406  & 86  & 37\% & 99.7\%  & 53.45  & 343.39 \\
10 & 144  & 773  & 9   & 13\% & 84.7\%  & 126.88 & 324.98 \\
08 & 1033 & 551  & 103 & 33\% & 79.7\%  & 648.49 & 3785.02 \\
\bottomrule
\end{tabular}
\end{table*}

Table~\ref{tab:flights} gives per-flight results. Seven flights localise with
median error between 8.4\,m and 24.8\,m and coverage at or above 99.7\%. Three
fail.

\subsection{Component comparison}
On flight 03, comparing the three approaches directly: visual odometry alone
covers every frame but drifts to 2803\,m; single-frame map matching achieves
5.46\,m median but only on the 76.6\% of frames it can solve, at 3.80 matching
calls per frame; the fused system covers 100\% of frames at 6.20\,m median with
1.52 calls per frame. The single-frame baseline produces no gross outliers
beyond 100\,m --- when retrieval proposes the wrong region, verification
rejects it and the system reports nothing rather than reporting a wrong
position.

\subsection{What separates working flights from failing ones}
\label{sec:law}
Performance varies by a factor of 78 across flights, which raises the question
of whether the system fails on some flights or whether some flights are
unmatched\-able. We measure the latter directly: with the true position given,
what fraction of frames match the basemap at all? This is a property of the
data, since no search is involved.

The answer separates the two groups. Flights with match rate $\geq 50\%$ give
median errors of 8.35--24.79\,m (median 15.44\,m) at coverage $\geq$99.7\%;
those below give 53--648\,m at 80--85\% coverage. The correlation between match
rate and log error is $-0.764$. Altitude is not the discriminator: the 2572\,m
flight works and the 551\,m flight fails.

This is practically useful because match rate needs no flight. Given a planned
route and a basemap, it can be measured in advance, converting deployability
from a question into a measurement.

\subsection{Ablation}
Removing components one at a time (300 frames, fixed seed): the boresight
correction is by far the largest single contributor (17.13\,m without it
against 6.62\,m with); removing odometry barely changes the median but
destroys the tail (90th percentile 368\,m against 16\,m), which is precisely
what a motion model is for; removing measurement-driven injection costs
1.0\,m median and 3.3\,m at the 90th percentile.

Two components do not earn their place and we report this rather than claiming
otherwise: particle count is nearly irrelevant (7.08\,m with 100 particles,
6.80\,m with 2000, 6.62\,m with 600), and removing the likelihood outlier floor
changes nothing measurable, because gating and injection already filter bad
fixes upstream.

\subsection{Robustness}
Across twenty conditions on 300 frames, the pattern is that the system is
indifferent to radiometry and sensitive to texture. Fog at maximum severity
costs 0.6\,m; occlusion covering a third of the frame costs 0.6\,m; resolution
loss costs 0.9\,m. Motion blur is the failure mode: moderate blur takes median
error from 6.62\,m to 58.31\,m, and heavy blur to 500\,m. Extreme darkness
prevents initialisation entirely, producing no position at all.

Discarding satellite fixes at random costs little until they become scarce:
6.62\,m at no dropout, 7.38\,m with half of all fixes discarded, 10.76\,m at
75\% and 41.55\,m at 90\%.

Blur can be partially mitigated by equalising the domain: since the basemap
tile is sharp and depicts the same ground at the same scale, the sharpness gap
between query and tile measures the blur, and blurring the tile to match
reduces median error under sustained blur from 44.47\,m to 20.00\,m and the
90th percentile from 214\,m to 91\,m. It does not restore clean performance;
heavy vibration remains a genuine limit.

\section{What Makes a Matcher Useful Here}
\label{sec:matcher}

Section~\ref{sec:law} identifies match quality as the binding constraint, which
suggests replacing LoFTR with a stronger matcher. RoMa~\cite{roma} leads
current matching benchmarks. Measured at known true positions it appears to
solve the problem outright: match rate on the three failing flights rises from
12--42\% to 96--100\%, within 2.7\,GB of VRAM.

End-to-end it is far worse: 1669\,m on flight 01, where LoFTR gives 22\,m.

The bake-off asked the wrong question. It showed each matcher only the
\emph{correct} tile, whereas a localization system spends most of its effort
deciding whether a candidate is the right place at all. Matching each frame
against a crop from a random, unrelated part of the map gives:

\begin{center}
\small
\begin{tabular}{lrrr}
\toprule
 & Correct place & Wrong place & Ratio \\
\midrule
LoFTR & 709 inliers & \textbf{0} & 709$\times$ \\
RoMa  & 4598 inliers & 341 & 13.5$\times$ \\
\bottomrule
\end{tabular}
\end{center}

LoFTR returns nothing on unrelated imagery; RoMa returns hundreds of
correspondences. Its apparent advantage was an artifact of a metric that only
measured the easy direction.

Recalibrating the acceptance threshold does not rescue it. Testing raw inlier
count, inlier ratio and matcher confidence on correct-versus-wrong pairs, LoFTR
admits a threshold with zero errors under two of the three criteria; for RoMa
no threshold exists under any criterion, as some wrong places outscore some
correct ones (best accuracy 98.1\%).

End-to-end, with thresholds tuned in RoMa's favour, flight 08 makes the cost
concrete. LoFTR produces a position on 4.5\% of frames, median error 582\,m ---
it is reporting that it does not know. RoMa produces one on 95.5\% of frames
and is 6250\,m wrong on average. In flight the first behaviour triggers a
handover to inertial navigation; the second flies the aircraft six kilometres
off course without indication.

We therefore keep LoFTR, not because it is the stronger matcher --- it is not
--- but because this task requires a property that matching benchmarks do not
measure. A matcher is useful for localization in proportion to its willingness
to return nothing.

\section{Limitations}

All ten flights come from one dataset and one country, and although altitudes
span 406--2572\,m and dates 2016--2023, the capture system and basemap class
are common to all. Three of the ten fail, and while the cause is measured, it
is also a real operational limit. Processing is offline at 866\,ms per frame on
a laptop GPU; frames arrive every 7\,s in this data, so it is comfortably
real-time for these flights, but it has not been run on embedded hardware.
Attitude and altitude are assumed available from non-GNSS sensors, and
Section~\ref{sec:selfcal} shows those sensors are themselves imperfect. The
planar homography assumption holds for near-nadir frames and degrades with
obliquity; building relief displacement is the dominant residual error source
and is why the achievable floor is near 6\,m rather than near 1\,m.

\section{Conclusion}

Treating GNSS-denied visual localization as sequential estimation rather than
per-frame retrieval changes which problem is hard. Retrieval accuracy is not
the constraint; the constraint is that many frames cannot be matched at all,
and a motion model covers exactly those. Across ten real flights the fused
system localises 99.7\% or more of frames at 8--25\,m median wherever the
basemap and the camera depict a recognisably similar world, and the fraction of
frames satisfying that condition can be measured before flight.

The same match geometry that yields position also yields the platform's own
heading error, allowing a magnetometer to be calibrated against the map with no
positioning information. And the matcher study suggests a criterion the
matching literature does not currently report: for localization, selectivity
dominates sensitivity. A matcher that always answers is worse than one that
sometimes declines.

Code, data preparation and every measurement reported here are available at
\url{https://github.com/YusufGUNEL/GeoAnchor}.

\begin{thebibliography}{00}
\bibitem{mgs2net} M. Li \textit{et al.}, ``(MGS)$^2$-Net: Macro/micro-geometric
structure and scale modeling for cross-view geo-localization,'' arXiv:2602.10704, 2026.
\bibitem{dinogfsa} B. Hu \textit{et al.}, ``DINO-GFSA: LoRA-adapted DINOv3 with
semantic gating for UAV self-positioning,'' arXiv:2606.00784, 2026.
\bibitem{offnadir} Q. Qiao \textit{et al.}, ``OffNadirLoc: A benchmark for large
off-nadir view geo-localization,'' arXiv:2607.19951, 2026.
\bibitem{relate} H. Jiang \textit{et al.}, ``ReLATE: Reliability-guided evidence
fusion for degraded UAV--satellite matching,'' arXiv:2607.25524, 2026.
\bibitem{scaleaware} Y. Ye \textit{et al.}, ``Scale-aware UAV-to-satellite
localization with semantic anchors,'' arXiv:2603.07535, 2026.
\bibitem{bevpatchpf} D. Lee \textit{et al.}, ``BEV-Patch-PF: Particle filtering
with learned aerial patch matching for GPS-free localization,'' arXiv:2512.15111, 2025.
\bibitem{downes} L. M. Downes \textit{et al.}, ``City-wide street-to-satellite
image geolocalization,'' arXiv:2203.05612, 2022.
\bibitem{shan} M. Shan \textit{et al.}, ``Google map aided visual navigation for
UAVs in GPS-denied environment,'' arXiv:1703.10125, 2017.
\bibitem{loftr} J. Sun \textit{et al.}, ``LoFTR: Detector-free local feature
matching with transformers,'' in \textit{CVPR}, 2021.
\bibitem{roma} J. Edstedt \textit{et al.}, ``RoMa: Robust dense feature
matching,'' in \textit{CVPR}, 2024.
\bibitem{dinov2} M. Oquab \textit{et al.}, ``DINOv2: Learning robust visual
features without supervision,'' \textit{TMLR}, 2024.
\bibitem{uavvisloc} W. Xu \textit{et al.}, ``UAV-VisLoc: A large-scale dataset
for UAV visual localization,'' arXiv:2405.11936, 2024.
\end{thebibliography}

\end{document}
